\subsection{\texorpdfstring{CANONICAL BASIS OF \(\mathfrak{sl}(2,k)\)}{CANONICAL BASIS OF \textbackslash mathfrak\{sl\}(2,k)}}

Lemma 1. Let \(\mathbf{A}\) be an associative algebra over \(k\) , \(H\) and \(X\) elements of \(\mathbf{A}\) such that \([H, X] = 2X\) .

\begin{enumerate}
\def\labelenumi{(\roman{enumi})}
\item
  \([H,X^n ] = 2nX^n\) for any integer \(n\geq 0\)
\item
  If \(Z\) is an element of \(A\) such that \([Z, X] = H\) , then, for any integer \(n > 0\) ,
\end{enumerate}

\[
[ Z, X ^ {n} ] = n X ^ {n - 1} (H + n - 1) = n (H - n + 1) X ^ {n - 1}.
\]

The map \(T \mapsto [H, T]\) from \(A\) to \(A\) is a derivation, which implies (i). With the assumptions in (ii),

\[
\begin{array}{l} [ Z, X ^ {n} ] = \sum_ {i + j = n - 1} X ^ {i} H X ^ {j} \\ \qquad = \sum_ {i + j = n - 1} (X ^ {i} X ^ {j} H + X ^ {i} 2 j X ^ {j}) \\ \qquad = n X ^ {n - 1} H + 2 X ^ {n - 1} \frac {n (n - 1)}{2} \\ \qquad = n X ^ {n - 1} (H + n - 1). \end{array}
\]

On the other hand, \(X^{n-1}(H+n-1)=(H-n+1)X^{n-1}\) by (i). Q.E.D.

Recall that we denote by \(\mathfrak{sl}(2,k)\) the Lie algebra consisting of the square matrices of order 2, trace zero, and with entries in k. This Lie algebra is simple of dimension 3 (Chap. I, §6, no. 7, Example). The canonical basis of \(\mathfrak{sl}(2,k)\) is the basis \((X_{+},X_{-},H)\) , where

\[
X _ {+} = \left( \begin{array}{c c} 0 & 1 \\ 0 & 0 \end{array} \right) \quad X _ {-} = \left( \begin{array}{c c} 0 & 0 \\ - 1 & 0 \end{array} \right) \quad H = \left( \begin{array}{c c} 1 & 0 \\ 0 & - 1 \end{array} \right).
\]

We have

\[
[ H, X _ {+} ] = 2 X _ {+} \quad [ H, X _ {-} ] = - 2 X _ {-} \quad [ X _ {+}, X _ {-} ] = - H.\tag{1}
\]

Since the identity representation of \(\mathfrak{sl}(2,k)\) is injective, H is a semi-simple element of \(\mathfrak{sl}(2,k)\) and \(X_{+}, X_{-}\) are nilpotent elements of \(\mathfrak{sl}(2,k)\) (Chap. I, §6, no. 3, Th. 3). By Chap. VII, §2, no. 1, Example 4, kH is a Cartan subalgebra of \(\mathfrak{sl}(2,k)\) . The map \(U \mapsto -^{t}U\) is an involutive automorphism of the Lie algebra \(\mathfrak{sl}(2,k)\) , called the canonical involution of \(\mathfrak{sl}(2,k)\) ; it transforms \((X_{+}, X_{-}, H)\) into \((X_{-}, X_{+}, -H)\) .

Lemma 2. In the enveloping algebra of \(\mathfrak{sl}(2,k)\) ,

\[
[ H, X _ {+} ^ {n} ] = 2 n X _ {+} ^ {n} \quad [ H, X _ {-} ^ {n} ] = - 2 n X _ {-} ^ {n}
\]

for any integer \(n \geq 0\) , and

\[
[ X _ {-}, X _ {+} ^ {n} ] = n X _ {+} ^ {n - 1} (H + n - 1) = n (H - n + 1) X _ {+} ^ {n - 1}
\]

\[
[ X _ {+}, X _ {-} ^ {n} ] = n X _ {-} ^ {n - 1} (- H + n - 1) = n (- H - n + 1) X _ {-} ^ {n - 1}
\]

if n \textgreater{} 0.

The first and third relations follow from Lemma 1. The others can be deduced from them by using the canonical involution of \(\mathfrak{sl}(2,k)\) .

